\section{Open questions}

Five open questions surfaced by this replication that the paper does not
close and that are testable at modest cost:

\begin{enumerate}

\item \textbf{Does nu-VQE generalize to open-quantum-system (OQS) VQE
(Lindbladian steady states)?}
The estimator $\text{Tr}[JHJ\rho]/\text{Tr}[J^2\rho]$ is well-defined for any
positive $\rho$, so nothing formally forbids the extension to
non-equilibrium steady states. But the paper only tests unitary-evolution +
noise; it does not check whether the diagonal Jastrow parameterization is
expressive enough to compensate for structured dissipation (e.g. non-Z-axis
dephasing).
\emph{Next steps:} run nu-VQE on a small dissipative Bose-Hubbard or
driven-dissipative TFIM against exact Lindbladian vectorization; sweep
dissipation strength; measure whether the $\sim$10--30$\times$ error
reduction survives.

\item \textbf{How does nu-VQE compose with modern imaginary-time-evolution
(ITE) variants (QITE, VarQITE, projected ITE)?}
Both nu-VQE and QITE effectively realize non-unitary transformations of the
ansatz; nu-VQE's Jastrow is diagonal (free at circuit depth), QITE
compiles a general non-unitary at real quantum cost. The paper does not
compare them.
\emph{Next steps:} on H$_2$/6-31G, benchmark plain VQE, nu-VQE, 1-step
VarQITE at matched depth, and nu-VQE seeded from a VarQITE-relaxed state.
Report error vs total circuit depth and vs classical parameter count.

\item \textbf{Is the noise-error-reduction ratio robust to non-depolarizing
noise channels (amplitude damping, coherent over-rotation, correlated
crosstalk)?}
Depolarizing noise commutes cleanly with diagonal $J$; amplitude damping
does not. Our replication used only depolarizing; the paper's ibmq
calibration was not decomposed to isolate contributions.
\emph{Next steps:} repeat nu-VQE vs VQE under four separately controlled
channels: (a) pure depolarizing, (b) T1-only amplitude damping,
(c) coherent $R_y$ over-rotation by fixed $\epsilon$, (d) correlated $ZZ$
crosstalk on the CNOT. If (b)--(d) show $<2\times$ reduction while (a)
shows $\sim 30\times$, the headline is conditional on channel type.

\item \textbf{Can an ML-selected / learned non-unitary structure outperform
the hand-picked pairwise-Z Jastrow, without inflating classical parameter
count?}
The paper picks $J = \exp(\text{1-body } Z + \text{2-body } ZZ)$ by
chemistry intuition. Optimality on the (accuracy, parameter-count) frontier
is not shown.
\emph{Next steps:} build a candidate library of $J$-generators up to
weight 3 (including $X$-basis and mixed strings), apply group-Lasso / $L_0$
regularization to select the smallest active set that hits chemical
accuracy under a fixed noise model, and publish the Pareto curve.

\item \textbf{Does nu-VQE extend to non-Hermitian effective Hamiltonians
(PT-symmetric systems, resonance / Siegert states, complex-scaled
molecular problems)?}
The Rayleigh-like quotient bound fails for non-Hermitian $H$, but the
non-unitary $J$ is a natural candidate for an approximate similarity
transform / biorthogonal partner.
\emph{Next steps:} on a small PT-symmetric 2-site Hatano-Nelson model,
compare (i) standard VQE minimizing $|\langle\psi|H|\psi\rangle|^2$,
(ii) nu-VQE with the paper's $J$, (iii) a biorthogonal ansatz where $J$
serves as an approximate similarity transform. Reproducing the
exceptional-point structure would open a new application domain beyond
ground-state chemistry.

\end{enumerate}
